\chapter{Cola de comandos}
\label{ap:cola}

\renewcommand{\arraystretch}{1.2}

\small
\subsection*{\color{verdeMB}Configuração}
\begin{tabularx}{\linewidth}{@{}>{\leavevmode\ttfamily\small\color{azulMB}}p{6.4cm} X@{}}
git config --global user.name "..." & seu nome \\
git config --global user.email "..." & seu e-mail \\
git config --list & ver a configuração \\
\end{tabularx}

\subsection*{\color{verdeMB}Começar}
\begin{tabularx}{\linewidth}{@{}>{\leavevmode\ttfamily\small\color{azulMB}}p{6.4cm} X@{}}
git init -b main & criar repositório \\
git clone <URL> & copiar um remoto \\
\end{tabularx}

\subsection*{\color{verdeMB}Registrar}
\begin{tabularx}{\linewidth}{@{}>{\leavevmode\ttfamily\small\color{azulMB}}p{6.4cm} X@{}}
git status & estado atual \\
git diff & o que mudou \\
git add arquivo & preparar \\
git commit -m "..." & registrar \\
git commit -am "..." & preparar e registrar os já acompanhados \\
git log --oneline --graph --all & histórico desenhado \\
\end{tabularx}

\subsection*{\color{verdeMB}Branches}
\begin{tabularx}{\linewidth}{@{}>{\leavevmode\ttfamily\small\color{azulMB}}p{6.4cm} X@{}}
git branch & listar \\
git switch -c nome & criar e entrar \\
git switch nome & trocar \\
git branch -m antigo novo & renomear \\
git merge nome & trazer \texttt{nome} para a atual \\
git merge --abort & cancelar o merge \\
\end{tabularx}

\subsection*{\color{verdeMB}Conflitos}
\begin{tabularx}{\linewidth}{@{}>{\leavevmode\ttfamily\small\color{azulMB}}p{6.4cm} X@{}}
git restore --ours arq & manter o lado da branch atual \\
git restore --theirs arq & manter o lado que chega \\
git add arq & marcar como resolvido \\
\end{tabularx}

\subsection*{\color{verdeMB}Remoto}
\begin{tabularx}{\linewidth}{@{}>{\leavevmode\ttfamily\small\color{azulMB}}p{6.4cm} X@{}}
git remote -v & listar remotos \\
git remote add origin <URL> & adicionar remoto \\
git fetch origin & baixar sem mesclar \\
git pull & baixar e mesclar \\
git push origin branch & enviar \\
\end{tabularx}

\subsection*{\color{verdeMB}Desfazer}
\begin{tabularx}{\linewidth}{@{}>{\leavevmode\ttfamily\small\color{azulMB}}p{6.4cm} X@{}}
git restore arq & descartar mudança não preparada \\
git restore --staged arq & tirar da preparação \\
git revert <commit> & desfazer com novo commit \\
\end{tabularx}

\subsection*{\color{verdeMB}GitLab CLI}
\begin{tabularx}{\linewidth}{@{}>{\leavevmode\ttfamily\small\color{azulMB}}p{6.4cm} X@{}}
glab auth login & autenticar \\
glab mr create & criar MR \\
glab mr list & listar MRs \\
glab mr view & ver o MR \\
glab mr merge & fazer o merge \\
glab ci status & ver o pipeline \\
\end{tabularx}


% -----------------------------------------------------------------------------
\chapter{Glossário}
\label{ap:glossario}

\begin{description}[style=nextline, font=\bfseries\color{azulMB}, leftmargin=0pt, itemsep=4pt]
\item[Agente de IA] Programa baseado em modelo de linguagem que executa tarefas num ambiente (ler arquivos, rodar comandos, escrever documentos), e não só conversa.
\item[Ancestral comum] Commit a partir do qual duas branches se separaram. Sem ele, o Git recusa o merge comum.
\item[Área de preparação (\emph{staging})] Onde ficam as mudanças escolhidas para o próximo commit.
\item[Bare repository] Repositório sem pasta de trabalho, só com o histórico. É o que um servidor guarda.
\item[Branch] Linha de desenvolvimento; tecnicamente, um ponteiro para um commit.
\item[Branch protegida] Branch que não aceita push direto; mudanças só por Merge Request.
\item[CI/CD] Integração e entrega contínuas: testes e implantação automáticos a cada mudança.
\item[Commit] Registro de um estado do projeto, com autor, data e mensagem.
\item[Conflito] Quando duas mudanças na mesma região de um arquivo não podem ser combinadas automaticamente.
\item[Fast-forward] Merge em que o Git só avança o ponteiro, porque não houve trabalho paralelo.
\item[GAP] Lacuna entre o estado atual (AS-IS) e o desejado (TO-BE), ou informação que falta e não deve ser inventada.
\item[HEAD] Referência para onde você está agora (em geral, a branch atual).
\item[Merge] Integração de duas linhas de histórico.
\item[Merge Request (MR)] Pedido para integrar uma branch em outra, com revisão. No GitHub: \emph{Pull Request}.
\item[\texttt{origin}] Nome padrão do repositório remoto principal.
\item[\texttt{ours} / \texttt{theirs}] Num merge: o lado da branch atual / o lado que está chegando.
\item[Pipeline] Sequência automática de etapas (compilar, testar, publicar) executada pelo GitLab.
\item[Prompt] Instrução dada a um modelo ou agente de IA.
\item[Remoto] Cópia do repositório em outro lugar, como o GitLab.
\item[Repositório] Pasta cujo histórico é controlado pelo Git.
\item[Self-Managed] GitLab instalado e administrado pela própria organização.
\item[Token] Na IA, pedaço de texto usado para medir custo e limite. No GitLab, credencial de acesso (trate como senha).
\end{description}

% -----------------------------------------------------------------------------
\chapter{Plano da apresentação}
\label{ap:plano}

Roteiro de tempo da capacitação de 01/10/2026, das 10h às 11h. Os números de slide se referem a \arq{apresentacao/apresentacao.pdf}.

\begin{center}
\renewcommand{\arraystretch}{1.4}
\begin{tabularx}{\linewidth}{@{}>{\bfseries\color{azulMB}}l c c >{\raggedright\arraybackslash}X@{}}
\toprule
\textbf{Horário} & \textbf{Duração} & \textbf{Slides} & \textbf{Conteúdo} \\
\midrule
10h00 & 3 min & 1--2 & Abertura, agenda e objetivos \\
10h03 & 19 min & 3--15 & Fundamentos: por que versionar, Git, GitLab, Self-Managed, branch, Merge Request, comandos, fluxo, conflito, \texttt{glab} \\
10h22 & 15 min & 16--25 & Problema 1: cenário, históricos sem relação, os oito passos, \texttt{ours} e \texttt{theirs} \\
10h37 & 18 min & 26--35 & Problema 2: IA como ferramenta, evolução dos prompts, falhas e regras, números, fases, travas, idioma, entregas \\
10h55 & 5 min & 36--39 & Encerramento, mensagem final e perguntas \\
\bottomrule
\end{tabularx}
\end{center}

\begin{dica}
Se o tempo apertar, os slides 8 (recursos do GitLab), 12 (tabela de comandos) e 35 (o que o prompt entrega) podem ser passados rapidamente: o conteúdo deles está completo nesta apostila.
\end{dica}

% -----------------------------------------------------------------------------
\chapter{Créditos e licenças}
\label{ap:creditos}

Todo o material de terceiros usado aqui está salvo localmente, na pasta \arq{assets/} (imagens e fontes), e o manual da marca em \arq{materials/manual-da-marca/}, para que a apostila e a apresentação possam ser recompiladas sem internet.

\begin{center}
\renewcommand{\arraystretch}{1.35}
\footnotesize
\begin{tabularx}{\linewidth}{@{}>{\raggedright\arraybackslash}p{3.4cm} >{\raggedright\arraybackslash}X >{\raggedright\arraybackslash}p{3.3cm}@{}}
\toprule
\textbf{Item} & \textbf{Origem} & \textbf{Licença} \\
\midrule
Manual de Identidade Visual da Marinha do Brasil (cores e tipografia) & CCEM, dez/2025 (\arq{materials/manual-da-marca/}) & Material institucional da Marinha do Brasil \\
Foto da Ilha das Cobras & Helder da Rocha, via Wikimedia Commons (\enquote{IlhaCobras.jpg}), reduzida & CC BY-SA 2.0 \\
Logotipo do Git & Jason Long, via Wikimedia Commons & CC BY 3.0 \\
Logotipo do GitLab (tanuki) & GitLab Inc., \arq{gitlab-com/gitlab-artwork} & CC BY-NC-SA 4.0; marca registrada da GitLab Inc. \\
Símbolo do GitHub & GitHub Octicons, via Wikimedia Commons & MIT; marca da GitHub Inc. \\
Símbolo do Claude & Wikimedia Commons & CC0; marca da Anthropic \\
Tirinha xkcd \#1597 \enquote{Git} & Randall Munroe, \url{https://xkcd.com/1597/} & CC BY-NC 2.5 \\
Fonte Montserrat & Julieta Ulanovsky e colaboradores & SIL Open Font License 1.1 \\
Fonte Bebas Neue & Dharma Type, via Google Fonts & SIL Open Font License 1.1 \\
Fonte Fira Mono & Mozilla / Carrois & SIL Open Font License 1.1 \\
\bottomrule
\end{tabularx}
\end{center}

As marcas de terceiros (Git, GitLab, GitHub, Claude) aparecem apenas para identificar as ferramentas de que o texto trata. As licenças NC (não comercial) são compatíveis com o uso deste material: capacitação interna, sem fins lucrativos.

\subsection*{Identidade visual}

As cores seguem o Manual de Identidade Visual da Marinha do Brasil: Azul Marinha \texttt{\#050F41}, Amarelo Ouro \texttt{\#FAB932} e Verde Brasil \texttt{\#079551}. A tipografia de apoio oficial é a Montserrat, usada no texto. As tipografias de título do manual (Octin College e Bebas Neue Pro) são comerciais. No lugar delas, usamos a Bebas Neue livre, do mesmo desenho da alternativa oficial.

% -----------------------------------------------------------------------------
\chapter{Para continuar estudando}
\label{ap:referencias}

\begin{itemize}
  \item \textbf{Pro Git}, de Scott Chacon e Ben Straub. O livro de referência do Git, gratuito e com tradução para o português: \url{https://git-scm.com/book/pt-br/v2}.
  \item \textbf{Documentação oficial do Git}: \url{https://git-scm.com/docs}.
  \item \textbf{Documentação do GitLab}: \url{https://docs.gitlab.com}. Veja em especial \emph{Merge requests} e \emph{Protected branches}.
  \item \textbf{GitLab CLI (\texttt{glab})}: \url{https://gitlab.com/gitlab-org/cli}.
  \item \textbf{Claude Code}: \url{https://docs.claude.com/en/docs/claude-code/overview}.
  \item \textbf{Learn Git Branching}, exercícios visuais e interativos de branch e merge, em português: \url{https://learngitbranching.js.org/?locale=pt_BR}.
\end{itemize}

\subsection*{Arquivos deste material}

\begin{center}
\renewcommand{\arraystretch}{1.3}
\small
\begin{tabularx}{\linewidth}{@{}>{\ttfamily\footnotesize}l X@{}}
\toprule
apresentacao/apresentacao.pdf & os slides da capacitação \\
apostila/apostila.pdf & esta apostila \\
model/SYSTEM\_DOCUMENTATION.md & o prompt final (V4) de documentação \\
roteiro\_versao\_01.md & o roteiro inicial, com os assuntos obrigatórios \\
roteiro\_versão\_final.md & o roteiro final \\
\bottomrule
\end{tabularx}
\end{center}
